\section{Individual Contribution Statements}
The team has confirmed that it contains six members.

\begin{description}[style=nextline,leftmargin=0pt]
\item[Member 1 -- Chuqiao Tang (2360692), Project Manager:] \pending{} contribution statement and Git evidence. The confirmed coordination role is Project Manager.
\item[Member 2 -- Jinyu Ma (2363433):] \pending{} contribution statement, project role, and Git evidence.
\item[Member 3 -- Muqing Xu (2362900):] My main responsibility was to reproduce the numerical experiments and verify that the figures used in the final report were generated correctly from the implemented methods and data. I checked the numerical workflow for the Robertson problem, including the fixed-step methods, reference-solution calculations, stability and convergence experiments, and the scripts used to generate the project figures.

I compared regenerated numerical outputs with the values reported in the presentation and report, and investigated discrepancies in quantities such as convergence orders, stability-step thresholds, reference-solution values, Newton-tolerance behaviour, and adaptive RK4 results. I also checked that figure labels, plotted quantities, logarithmic scales, and solver classifications matched the corresponding numerical data rather than being manually adjusted. Where results differed between versions, I traced the difference to changes in parameters, fitting ranges, solver settings, or diagnostic definitions before the final figure was accepted.

I also supported reproducibility by confirming which scripts and numerical outputs produced each figure and by checking that regenerated plots were consistent with the final reported results. My contribution focused on code reproduction, numerical verification, and figure validation rather than report writing.
\item[Member 4 -- Shuoran Fu (2362626):]My primary responsibility in this project was the writing, integration, and revision of the final report. I was responsible for organising the report in LaTeX and integrating the team’s mathematical derivations, numerical results, figures, and experimental findings into a coherent final document. I drafted and revised the written discussion across the main sections, including the Introduction, Numerical Methods, Mathematical and Stability Analysis, Reference Solution and Validation, Numerical Experiments, and Conclusion.  

I also revised the report systematically against the assessment rubric provided by the instructor. In particular, I checked whether the report clearly explained the mathematical methods, connected the stability analysis to the Robertson problem, interpreted the numerical experiments rather than only describing them, included matched-accuracy cost comparisons, and discussed limitations and conclusions appropriately.  

In addition, I maintained the report structure in the GitHub repository, including the section files under report/sections/, the abstract, references, and individual contribution statement. I standardised notation, equations, formatting, figure placement, captions, cross-references, and tables, and regularly updated the LaTeX source and compiled PDF.
\item[Member 5 -- Taiqi Zhang (2363684):] \pending{} contribution statement, project role, and Git evidence.
\item[Member 6 -- Yujia Ding (2362094):] Within our six-member Group C project investigating the stiff Robertson chemical-kinetics ODE problem, my primary responsibilities covered code implementation, code verification, and all figure reproduction for the presentation. I carried out code inspection and end-to-end validation for the full project codebase, checking compilation and reproducibility to ensure that the numerical routines for explicit Euler, RK4, and implicit Euler ran correctly without implementation bugs.

Based on verified simulation outputs, I reproduced every plot shown in the group slides, including stiffness-ratio curves, stability-region diagrams, four-track diagnostic plots, and work--precision cost graphs. Besides coding and visualisation work, I participated in revising and polishing the group report, refining the argument logic and checking numerical consistency between textual statements and plotted data.

I also joined regular group discussions and helped cross-check quantitative results such as the stable step-size bounds and the reference value $y_2(40)$. This work improved my practical understanding of stiff ODEs, numerical stability, and rigorous result validation.
\end{description}
